% GENERATED FILE. Edit canonical lessons, then run npm run generate:books.
% canonical-source-hash: fnv1a64:fe163e95534a8b6b
% canonical-lessons: SA-C27-so-be-it, SA-C27-well-done, SA-C27-enough, SA-C27-true, SA-C27-thus

\chapter{Answering With More Than Yes}
\label{ch:sa-replies}
\hlchaptermodality{27}

\begin{chapteropening}
\textbf{By the end of this chapter:} \emph{I can answer in Sanskrit with more than yes and no, and hear the be-verb sitting inside three of the answers.}

Complete SA-C27-thus and demonstrate the chapter promise.
\end{chapteropening}

\section[astu]{\sk{अस्तु} (astu) — so be it, all right}

\label{lesson:SA-C27-so-be-it}



\begin{quote}
Before the new word: what did \emph{vaidyaḥ} mean?
\end{quote}

\subsection*{You'll want to know: \sk{अस्तु}}
\textbf{\sk{अस्तु}} (\emph{astu}) — \textquotedblleft{}so be it, all right\textquotedblright{}.

You have had \textbf{\sk{आम्}} and \textbf{\sk{न}} for yes and no since the first chapter, and they are blunt. \textbf{\sk{अस्तु}} is the agreeable middle: \emph{let it be so}, \emph{all right}, \emph{fine}.

It is not a new word so much as a new face on one you know. \textbf{\sk{अस्ति}} (\emph{asti}), \textquotedblleft{}is\textquotedblright{}, was the third person of the root \emph{as}, \textquotedblleft{}to be\textquotedblright{}. \sk{अस्तु} is the same verb in the commanding form: \textbf{let it be}. Sanskrit answers a request by ordering the world to comply.

\begin{grammarlens}[title={What you've built}]
A reply softer than yes, built out of the be-verb you already had.
\end{grammarlens}

\subsection*{Your turn}
Say these aloud:

\begin{itemize}
\raggedright
  \item \emph{astu}
  \item it once more, slowly
  \item \emph{asti}, then \emph{astu}, and say what changed between them
\end{itemize}

\begin{tcolorbox}[breakable,colback=teal!4,colframe=teal!35!black,title={Before you move on}]
What does \sk{अस्तु} mean? (\textquotedblleft{}so be it, all right\textquotedblright{}.) Say it once more.
\end{tcolorbox}
\section[sādhu]{\sk{साधु} (sādhu) — well done, good}

\label{lesson:SA-C27-well-done}



\begin{quote}
Before the new word: what did \emph{astu} mean, and which verb was it built on?
\end{quote}

\subsection*{You'll want to know: \sk{साधु}}
\textbf{\sk{साधु}} (\emph{sādhu}) — \textquotedblleft{}well done, good\textquotedblright{}.

\textbf{\sk{साधु}} is what you call out when something has been done well — \emph{good}, \emph{well said}, the Sanskrit \emph{bravo}. An audience says it aloud at a concert.

Its root means \textquotedblleft{}to accomplish, to reach the goal\textquotedblright{}, so the word praises a thing for having arrived where it was going. The same word, standing on its own as a noun, names a holy man: \emph{a sādhu} is someone who has accomplished. English borrowed that sense and left the everyday one behind, exactly as it did with \sk{गुरुः}.

\begin{grammarlens}[title={What you've built}]
A second borrowed word restored to its ordinary use.
\end{grammarlens}

\subsection*{Your turn}
Say these aloud:

\begin{itemize}
\raggedright
  \item \emph{sādhu}
  \item it once more, slowly
  \item \emph{sādhu}, then \emph{guruḥ}, and say what English kept of each
\end{itemize}

\begin{tcolorbox}[breakable,colback=teal!4,colframe=teal!35!black,title={Before you move on}]
What does \sk{साधु} mean? (\textquotedblleft{}well done, good\textquotedblright{}.) Say it once more.
\end{tcolorbox}
\section[alam]{\sk{अलम्} (alam) — enough}

\label{lesson:SA-C27-enough}



\begin{quote}
Before the new word: what did \emph{sādhu} mean?
\end{quote}

\subsection*{You'll want to know: \sk{अलम्}}
\textbf{\sk{अलम्}} (\emph{alam}) — \textquotedblleft{}enough\textquotedblright{}.

\textbf{\sk{अलम्}} is \emph{enough}. Enough food, enough talking, enough of that. It is short, common, and does a lot of work for its size.

Set beside it a small word you have not been given yet and will hear constantly — \emph{alam} can also mean \emph{able}, \emph{equal to the task}. So the same two syllables cover \textquotedblleft{}that will do\textquotedblright{} and \textquotedblleft{}that is up to it\textquotedblright{}, which is roughly the distance English travels between \emph{enough} and \emph{sufficient}. No secure cousin outside India, so this one sits with \sk{केशः}.

\begin{grammarlens}[title={What you've built}]
A one-word answer, and a second word with no travel history.
\end{grammarlens}

\subsection*{Your turn}
Say these aloud:

\begin{itemize}
\raggedright
  \item \emph{alam}
  \item it once more, slowly
  \item \emph{alam}, then \emph{keśaḥ}, and name what those two have in common
\end{itemize}

\begin{tcolorbox}[breakable,colback=teal!4,colframe=teal!35!black,title={Before you move on}]
What does \sk{अलम्} mean? (\textquotedblleft{}enough\textquotedblright{}.) Say it once more.
\end{tcolorbox}
\section[satyam]{\sk{सत्यम्} (satyam) — true; the truth}

\label{lesson:SA-C27-true}



\begin{quote}
Before the new word: what did \emph{alam} mean, and what did \emph{astu} mean?
\end{quote}

\subsection*{You'll want to know: \sk{सत्यम्}}
\textbf{\sk{सत्यम्}} (\emph{satyam}) — \textquotedblleft{}true; the truth\textquotedblright{}.

\textbf{\sk{सत्यम्}} answers a claim: \emph{true}. As a noun it is \emph{the truth}, and it is the word on the national emblem of India — \emph{satyam eva jayate}, \textquotedblleft{}truth alone conquers\textquotedblright{}.

Look at what it is made of. \emph{Sat-} is the being-form of that same root \emph{as} you met in \sk{अस्ति} and \sk{अस्तु}: \emph{what is}. Add an ending and truth is simply \textbf{the state of being what is}. English has a cousin, well hidden: \emph{sooth}, as in \emph{soothsayer} — a truth-sayer — and \emph{forsooth}.

\begin{grammarlens}[title={What you've built}]
A third word from the be-verb, and truth defined as what is.
\end{grammarlens}

\subsection*{Your turn}
Say these aloud:

\begin{itemize}
\raggedright
  \item \emph{satyam}
  \item it once more, slowly
  \item \emph{asti}, \emph{astu}, \emph{satyam} — three from one root, in that order
\end{itemize}

\begin{tcolorbox}[breakable,colback=teal!4,colframe=teal!35!black,title={Before you move on}]
What does \sk{सत्यम्} mean? (\textquotedblleft{}true; the truth\textquotedblright{}.) Say it once more.
\end{tcolorbox}
\section[evam]{\sk{एवम्} (evam) — thus, so}

\label{lesson:SA-C27-thus}



\begin{quote}
Before the last word of this run: what did \emph{satyam} mean, and what did \emph{sādhu} mean?
\end{quote}

\subsection*{You'll want to know: \sk{एवम्}}
\textbf{\sk{एवम्}} (\emph{evam}) — \textquotedblleft{}thus, so\textquotedblright{}.

\textbf{\sk{एवम्}} means \emph{thus}, \emph{so}, \emph{in this way}. On its own it answers a question: \emph{just so}.

Its real value is that it combines. Put it in front of the first word of this run and you get \textbf{\sk{एवम्} \sk{अस्तु}} (\emph{evam astu}) — \textquotedblleft{}so be it\textquotedblright{}, the whole formula, made of two words you now own separately. That is what a small vocabulary buys: not five answers, but the joins between them.

\begin{grammarlens}[title={What you've built}]
Five replies past yes and no: so be it, well done, enough, true, thus.
\end{grammarlens}

\subsection*{Your turn}
Say these aloud:

\begin{itemize}
\raggedright
  \item \emph{evam}
  \item it once more, slowly
  \item all five in order, then \emph{evam astu} as one phrase
\end{itemize}

\begin{tcolorbox}[breakable,colback=teal!4,colframe=teal!35!black,title={Before you move on}]
What does \sk{एवम्} mean? (\textquotedblleft{}thus, so\textquotedblright{}.) Say it once more.
\end{tcolorbox}
